\section{本章小结}
\label{sec:pm-summary}

本章从一个问题开始：允许迭代点暂时违反约束，用什么信号把它拉回来。

惩罚方法给违规计价。二次罚光滑，但边界上没有回推的力，有限的$\rho$总留下违规，$\rho\to\infty$时法向曲率随之增长，子问题越来越病态。一次罚在边界上有尖角，$\rho$超过最优乘子的最大绝对分量后最小点就是原问题的解，代价是目标不可微。“精确”的含义是解的对应，全局结论需要强对偶和对偶解存在，局部结论需要误差界或二阶条件。

对偶上升直接调整价格。解完Lagrange子问题，约束残差就是对偶函数的超梯度，目标可分时各块可以独立求解。它改进的是下界，迭代点通常不可行，步长受限，$q$不可微时单步未必上升；要得到原始解，需要平均、选解或修复，并且同时报告可行性和间隙。

增广Lagrange法把价格和平方罚合在一起。乘子提供大部分的力，罚项修正余下的部分，所以有限的$\rho$就够了。从对偶问题看，它是隐式的邻近点步，比显式的对偶上升稳定得多，凸问题中任何固定的$\rho>0$都收敛。代价是平方项把变量耦合起来，子问题更难。不等式要用正部形式，乘子要保持非负。

ADMM把两块变量各更新一次，用一次乘子更新代替联合求解。这样既保留乘子的稳定性，又能分别利用两块结构，代价是出现了对偶残差。原始残差和对偶残差要一起看，$\rho$改变时缩放乘子要同步换算。两块闭凸目标在鞍点存在时有收敛保证，多块或非凸的情形不能沿用。

\begin{BookInlineTable}
  \centering
  \small
  \begin{tabular}{@{}lll@{}}
    \toprule
    方法 & 主要条件 & 能保证什么\\
    \midrule
    二次罚 & 各轮全局解，$\rho\to\infty$ & 聚点全局最优\\
    一次罚 & 强对偶，$\rho$超过门槛 & 最小点集合相同\\
    对偶上升 & 残差来自全局最小点 & 超梯度；凸时平均点渐近可行\\
    ALM & 凸，强对偶，固定$\rho$ & 残差趋零，目标值收敛\\
    ADMM & 两块闭凸，鞍点存在 & 两种残差趋零，目标值收敛\\
    \bottomrule
  \end{tabular}
  \caption{本章各方法的主要条件与保证。局部结论和近似求解下的结论见各节。}
  \label{tab:pm-summary}
\end{BookInlineTable}

几种方法结束时要做的事是一样的：回到原约束检查可行性，回到Lagrange函数检查驻定，有不等式时检查互补；需要全局结论时，再用对偶值或者凸性给出证书。第~\ref{chap:sequential-quadratic-programming}章接着处理光滑非线性约束，用局部二次模型同时更新原始变量和乘子。
